% page 421 of the facsimile (image p-420 of the scan)
% block 165.1 x 246.3 mm ; left margin 36.9 mm
\pgc{15.240mm}{0.000mm}{
\l{\cel{50}413}
\l{}
\l{\sou{ostracismo}. (en\cel{1}\sou{Grekia antiqu}a)\cel{2}Exilo, dek-yar-dura, di civi-}
\l{\cel{3}tano qua divenis suspektata pro lua autoritato, lua povo, en}
\l{\cel{3}la civito. - DEFIRS.}
\l{}
\l{\sou{ostrogoto}. Persono qua ne konformigas sua konduto a la kustumi,}
\l{\cel{3}a la deci, qui regnas en la medio en qua lu vivas. -}
\l{}
\l{\sou{-ot-}\cel{1}.\cel{2}(\sou{gram.})\cel{2}Verbo-dezinenco di la participo di futureso,}
\l{\cel{3}pasiva, -}
\l{}
\l{\sou{otalgio}. (\sou{patol.)}\cel{2}Nevralgio di la orelo. -}
\l{}
\l{\sou{otario.}\cel{1}(\sou{zool.)}\cel{3}Speco di mar-hundo di qua la oreli salias.}
\l{\cel{3}- L.\cel{1}\sou{otaria.}\cel{2}EFI.}
\l{}
\l{\sou{otito}. (\sou{patol.}) Inflamuro di la orelo.}
\l{}
\l{\sou{otolito.(biol.)}\cel{2}Konkreciono tre mikra, kalkoza, qua existas}
\l{\cel{3}en la aud-organaro di animali diversa. -}
\l{}
\l{\sou{otomano}. Granda sidilo, sen dors-apogilo, sur qua onu povas}
\l{\cel{3}repozar segun la posturo di la Orientani.- DEFIRS.}
\l{}
\l{\sou{otuso.}\cel{2}(\sou{zool.)}\cel{2}Raptucelo noktala. - L.\cel{1}\sou{otus}.\cel{2}L.}
\l{}
\l{\sou{ovo}. I. Maso globatra, quan pondas la ucelino, jermo di la ucelo}
\l{\cel{3}futura, cirkondata do liquidi destinata nutror lu dum\cel{2}kelka}
\l{\cel{3}tempo, ed envelopata da shelo kalkoza. - II. Jermo-ovulo}
\l{\cel{3}che la animali vivipara qua, pos fekundigeso e developeso,}
\l{\cel{3}pasas de la ovario aden la utero. III. To quo havas la formo}
\l{\cel{3}di ovo. - IV. (\sou{tekn.)}\cel{2}Ornamento di arkitekturo, di juveli,}
\l{\cel{3}e c. qua havas la formo di ovo. - DEFIRS.}
\l{}
\l{\sou{ovacionar}. (\sou{trans.}) I. (en\cel{1}\sou{Roma, antique}). Ceremonio, min so-}
\l{\cel{3}lena kam triumfo, dum qua onu imolis mutonino, e la generalo}
\l{\cel{3}vinkinta avancis kavalkante sur kavalo. - II. Honoro agita}
\l{\cel{3}a personego da la procesionani, eskortani. - DEFIRS.}
\l{}
\l{\sou{ovario}. I. (\sou{zool.)}\cel{3}Organo di la femino, qua kontenas la ovuli}
\l{\cel{3}destinita produktar la feto, che la vivipari - la ovo, che}
\l{\cel{3}la ovipari. - II.(\sou{bot.}) Parto di la pistilo, qua kontenas}
\l{\cel{3}la ovuli destinita produktar la semini. - DEFIS.}
\l{}
\l{\sou{ovipara}. Qua produktas e fekundigas ovi de qui ekiros la yuni.}
\l{\cel{3}DEFIS.}
\l{}
\l{\sou{ovogonio}. (\sou{biol}.)Celuli qui rezultas de la genito di la elementi}
\l{\cel{3}jerma di la ovisako. -}
\l{}
\l{\sou{ovoida}. Di qua la formo aspektas quale ovo.}
\l{}
\l{\sou{ovu}lo. I. (\sou{zool.})\cel{2}Produkturo de la ovario, qua, pos fekindigeso,}
\l{\cel{3}divenas la feto che la vivipari - e la ovo, che la ovipari.}
\l{\cel{3}- II. (bot.) Produkturo de la ovario qua, pos fekundigeso, di-}
\l{\cel{3}venos la semino. - EFIS.}
}
